% Requires: amsmath, amssymb, booktabs.  Compact appendix for the thesis; full derivation: docs/technical_note.tex
\newcommand{\ii}{\mathrm{i}}

\chapter{Interface transmission: boundary-value problem and checks}
\label{app:interface_derivation}

This appendix gives the boundary-value problem solved in Sec.~\ref{sec:interface_AMM}. The
field expressions are standard (Auld, Ref.~\cite{Auld1973}; Every \emph{et al.},
Ref.~\cite{Every1984}); the complete derivation and the code are available at
\todo{repository URL / DOI}.

\section{Fields and coefficient matrix}

The crystal ($c$) occupies $z<0$ and the metal ($m$) $z>0$; all fields depend on $x$ and $z$
only, with time dependence $e^{-\ii\omega t}$. Write the displacement as
$\mathbf u=\nabla\phi+\nabla\times(\psi\hat{\mathbf y})$ with plane-wave potentials
$\phi=A\,e^{\ii\omega(px+sk_Lz-t)}$ and $\psi=B\,e^{\ii\omega(px+sk_Tz-t)}$, where
$s=\pm1$ is the sign of the vertical propagation direction. The horizontal slowness $p$ is the
same for every wave (Snell's law) and the vertical slowness is
$k_j=\sqrt{V_j^{-2}-p^2}$, taken with $\operatorname{Im}k_j\ge0$ so that evanescent waves
decay away from the interface. With the reduced fields $\tilde u=u/\omega$ and
$\tilde\sigma=\sigma/\omega^2$, the displacement and stress components per unit potential
amplitude are given in Table~\ref{tab:fields}, where $D=\rho-2\mu p^2$.

\begin{table}[H]
\centering
\caption{Reduced field vectors per unit potential amplitude, evaluated with the density and
shear modulus of the medium containing the wave.}
\label{tab:fields}
\begin{tabular}{lcccc}
\toprule
wave & $\tilde u_x$ & $\tilde u_z$ & $\tilde\sigma_{xz}$ & $\tilde\sigma_{zz}$\\
\midrule
L  & $\ii p$ & $\ii s k_L$ & $-2\mu s\,p\,k_L$ & $-D$\\
SV & $-\ii s k_T$ & $\ii p$ & $D$ & $-2\mu s\,p\,k_T$\\
\bottomrule
\end{tabular}
\end{table}

An incident wave ($s=+1$, normalised to unit displacement amplitude, $A_{\rm inc}=V_{\rm inc}$)
produces a reflected L and SV wave in the crystal ($s=-1$) and a transmitted L and SV wave
in the metal ($s=+1$), with unknown potential amplitudes
$\mathbf x=(A_r,B_r,A_t,B_t)^{\mathsf T}$. Continuity of traction and the spring law,
$\sigma^c_{xz}=\sigma^m_{xz}=K_\parallel(u_x^m-u_x^c)$ and
$\sigma^c_{zz}=\sigma^m_{zz}=K_\perp(u_z^m-u_z^c)$, give $M\mathbf x=\mathbf b$ with
\begin{equation}
  M=\begin{pmatrix}
  \tilde\sigma_{xz}^{rL} & \tilde\sigma_{xz}^{rT} & -\tilde\sigma_{xz}^{tL} & -\tilde\sigma_{xz}^{tT}\\[2pt]
  \tilde\sigma_{zz}^{rL} & \tilde\sigma_{zz}^{rT} & -\tilde\sigma_{zz}^{tL} & -\tilde\sigma_{zz}^{tT}\\[2pt]
  \tilde\sigma_{xz}^{rL}+\tilde K_\parallel\tilde u_x^{rL} & \tilde\sigma_{xz}^{rT}+\tilde K_\parallel\tilde u_x^{rT} &
   -\tilde K_\parallel\tilde u_x^{tL} & -\tilde K_\parallel\tilde u_x^{tT}\\[2pt]
  \tilde\sigma_{zz}^{rL}+\tilde K_\perp\tilde u_z^{rL} & \tilde\sigma_{zz}^{rT}+\tilde K_\perp\tilde u_z^{rT} &
   -\tilde K_\perp\tilde u_z^{tL} & -\tilde K_\perp\tilde u_z^{tT}
  \end{pmatrix},\quad
  \mathbf b=-\begin{pmatrix}
  \tilde\sigma_{xz}^{\rm inc}\\ \tilde\sigma_{zz}^{\rm inc}\\
  \tilde\sigma_{xz}^{\rm inc}+\tilde K_\parallel\tilde u_x^{\rm inc}\\
  \tilde\sigma_{zz}^{\rm inc}+\tilde K_\perp\tilde u_z^{\rm inc}
  \end{pmatrix},
\end{equation}
where $\tilde K_\parallel=\alpha_\parallel\mu_c/V_{T,c}$ and
$\tilde K_\perp=\alpha_\perp(\lambda_c+2\mu_c)/V_{L,c}$ follow from Eq.~\eqref{eq:alpha_omega}
(superscripts $r$, $t$, inc denote the reflected, transmitted and incident waves, and L, T
their type). The welded limit replaces rows 3--4 by continuity of $\tilde u_x$ and $\tilde u_z$;
the free-surface limit is $\tilde K\to0$. The matrix has no useful closed form for general
$\theta_i$, so it is inverted numerically at each angle. The amplitude coefficients are
$\Gamma_j=A_j/V_j$ (L) or $B_j/V_j$ (T), and the power coefficients are given by
Eq.~\eqref{eq:power_coeff}.

\section{SH channel}

For SH polarisation the same procedure gives a $2\times2$ system with the solution
\begin{equation}
  T_r=\frac{2}{1+b/a-\ii\,b/\tilde K_{SH}},\qquad a=\mu_ck_{T,c},\ \ b=\mu_mk_{T,m},\ \
  \tilde K_{SH}=\alpha_{SH}\frac{\mu_c}{V_{T,c}},
\end{equation}
with $\alpha_{SH}=\alpha_\parallel$ assumed.

\section{Energy conservation and validation}

For real $K$ the spring stores but does not dissipate energy, so the four power fractions of
Eq.~\eqref{eq:power_coeff} sum to one at every angle; the numerical solution satisfies this to
$10^{-8}$ for the welded, spring and free-surface cases and for SH. The implementation was
also checked against three closed-form limits:
\begin{itemize}
\item \emph{normal incidence}, where the transmitted power is
\begin{equation}
  W_{\rm trans}=\frac{4r}{(1+r)^2+r^2/\alpha^2},\qquad r=\frac{Z_m}{Z_c},\ \ Z=\rho V,
  \label{eq:normal_spring}
\end{equation}
reducing to $4r/(1+r)^2$ for $\alpha\to\infty$ (for sapphire/Ta and $\alpha=1$ it gives $0.739$);
\item the \emph{SH solution} above, for $\theta_i=10^\circ$, $40^\circ$ and $70^\circ$;
\item the \emph{free-surface} amplitudes of an incident SV wave,
$\Gamma_{T\to L}=-\kappa\sin4\theta/(\sin2\theta\sin2\theta_L+\kappa^2\cos^22\theta)$ with
$\kappa=V_L/V_T$ and $\sin\theta_L=\kappa\sin\theta$, which also gives the peak
$|\Gamma_{T\to L}|^2=14.7$ at $\theta_c=36.4^\circ$.
\end{itemize}
All agree with the numerical solver to $10^{-9}$ or better (the $10^{-6}$ for the
normal-incidence limit reflects the finite angle used to approximate $\theta\to0$). In addition, the welded flux-averaged
transmission of the Al/sapphire and Nb/sapphire interfaces computed with the velocities of Kaplan~\cite{Kaplan1979}
agrees with his tabulated values to within $1.3\%$.
